\section{Background}
\label{sec:related}

\subsection{Existing Selectivity Metrics and Their Comparison}
\label{sec:related:metrics}

Several metrics have been proposed to quantify kinase inhibitor selectivity
from profiling panel data.
Karaman et al.\cite{Karaman2008} introduced the selectivity score (S-score),
defined as the fraction of panel kinases inhibited above a fixed concentration
threshold.
Uitdehaag and Zaman\cite{Uitdehaag2011} proposed the selectivity entropy,
which treats the normalized binding profile as a probability distribution and
computes its Shannon entropy. Low entropy indicates activity concentrated on
few targets.
Graczyk\cite{Graczyk2010} adapted the Gini coefficient from economics, treating
the distribution of a compound's activity across the panel as an inequality
distribution. A value near 1 indicates activity concentrated on few kinases and
near 0 indicates activity spread evenly.
Klaeger et al.\cite{Klaeger2017} introduced the concentration- and
target-dependent selectivity (CATDS) score, which quantifies the reduction in
binding to a given target relative to the summed reduction across all detected
targets at a specified inhibitor concentration (see Discussion for why we
did not implement it as a fifth measure).
Bosc et al.\cite{Bosc2017} subsequently proposed two further metrics, the
window score and ranking score, and compared them against four previously
published metrics (the S-score, Gini coefficient, selectivity entropy, and a
partition index) on the Davis and Anastassiadis datasets, showing that
compound rankings differ substantially depending on the metric applied and that
certain compounds cannot be scored under some definitions when threshold
conditions are not satisfied. For example, ratio-based scores are undefined
for a compound with no kinase above the activity threshold, since there is then
no primary target or off-target to compare.

Reconciliation attempts have remained limited in scope.
Bajorath et al.\cite{Bajorath2018} compared CATDS-based selectivity
assessments from Klaeger et al.\ with activity data mined from ChEMBL for the
same compounds, finding broad consistency between the two sources.
The KInhibition platform\cite{KInhibition2018} explicitly acknowledged that
multiple existing metrics yield conflicting results for the same compound, but
addressed this by introducing its own composite score (combining on-target
inhibition with an off-target inhibition penalty) rather than reconciling the
disagreement among the existing definitions.
More recently, the same group's built a KIRHub resource\cite{Gujral2026}
on a
new profiling campaign of 86 approved kinase inhibitors against 758 kinases
and disease-associated mutant variants. This resource offers users a choice 
between a further
named metric (the KISS score, an on-target/off-target balance) and a separate
Gini-based score, thus proliferating named, ad hoc definitions.

Two further notable gaps surface in existing work. The first is that
is that each metric is evaluated at fixed
parameter values. Thus, the sensitivity of rankings to parameter choice has not
been systematically characterized. The second i that
the structural properties of binding profiles that determine when definitions
agree or disagree have not been identified.
The present work addresses both gaps.

\subsection{Axiomatic Approaches to Ranking and Measurement}
\label{sec:related:axiomatic}

Precedent for an
axiomatic characterization of measurement functions is established.
Shannon\cite{Shannon1948} demonstrated that the entropy of a probability
distribution is uniquely determined, up to a multiplicative constant, by four
requirements: continuity, maximality on the uniform distribution,
expansibility, and additivity over independent subsystems.
Khinchin\cite{Khinchin1957} subsequently strengthened this uniqueness result.
The selectivity entropy of Uitdehaag and Zaman applies the Shannon formula to
normalized binding profiles but provides no analogous axiomatic justification.
The axioms that single out Shannon entropy, such as maximality on the uniform
distribution and additivity over independent subsystems, characterize a measure
of uncertainty rather than a measure of selectivity.

Arrow\cite{Arrow1951} established that no social welfare function satisfies
simultaneously the requirements of unrestricted domain, Pareto efficiency,
independence of irrelevant alternatives, and non-dictatorship.
The theorem is relevant here in two respects: it demonstrates that axiomatic
analysis of a ranking concept can yield impossibility,
and it shows that competing ranking functions reflect genuine tradeoffs between
properties that cannot all be achieved simultaneously.
Singleton and Booth\cite{Singleton2022} applied this methodology to truth
discovery, i.e. the problem of identifying true facts from conflicting reports of
unknown reliability, posing it as a social choice problem, formulating
candidate axioms, and proving that a subset cannot be jointly satisfied.

To our knowledge, this axiomatic style of analysis has not previously been
applied to any selectivity or promiscuity measure in drug discovery.
The present work proposes empirically motivated properties as a
foundation for such an analysis.
